\begin{frame}[fragile]{Global Value Numbering}
  Global Value Numbering (GVN) is an optimization that finds equivalent expressions and then picks a representative one, \textbf{removing} redundant computation
    \begin{columns}[c]
        \begin{column}{0.44\textwidth}
           \begin{lstlisting}[language=C++, frame=single, xleftmargin=4pt, xrightmargin=4pt]
t = x + 1;
if (t == y) {
    w = x + 1;
    foo(w);
}
           \end{lstlisting} 
        \end{column}
        \begin{column}{0.08\textwidth}
            \transform        
        \end{column}
        \begin{column}{0.44\textwidth}
           \begin{lstlisting}[language=C++, frame=single, xleftmargin=4pt, xrightmargin=4pt]
t = x + 1;
if (t == (*@\textcolor{UnipiRed}{y}@*)) {
    foo((*@\textcolor{UnipiRed}{y}@*));
}
           \end{lstlisting} 
        \end{column}
    \end{columns}
    This is legal if branching on poison is considered \textbf{immediate UB}, the transformation would be unsound otherwise
    \pdfpcnote{
GVN eliminates redundant computations by recognizing equivalent expressions.
Inside if (t == y), the compiler knows t and y are equal, so it replaces w = x+1
with y. This is only valid if the equality is a reliable fact — which requires
that branching on poison IS immediate UB, so the condition is guaranteed to hold
in the branch body.
}
\end{frame}
